% =============================================================================
% Section 1: Introduction
% =============================================================================
\section{Introduction}
\label{sec:introduction}

% ---------------------------------------------------------------------------
% 1.1 Clinical Motivation
% ---------------------------------------------------------------------------
\subsection{Clinical Motivation}
\label{sec:intro:clinical}

Meningiomas are the most common primary intracranial tumours, accounting for
approximately 39\% of all central nervous system neoplasms, with an
age-adjusted incidence rate of 9.51 per 100,000 population
\cite{ostrom2023brain}. The majority are classified as WHO Grade~I (benign),
and a substantial fraction are discovered incidentally on neuroimaging
performed for unrelated indications \cite{vernooij2007incidental}. For these
incidentally discovered meningiomas, clinical management hinges on a
deceptively simple question: \emph{will this tumour grow?}

Current clinical practice relies on serial MRI surveillance with manual
volumetric assessment at intervals of 3--12 months
\cite{hashimoto2012natural,ius2020management}. This approach has three
fundamental limitations. First, manual volumetry is labour-intensive and
subject to considerable inter-observer variability, with reported coefficient
of variation exceeding 15\% for small tumours \cite{defined2016interobserver}.
Second, the decision to intervene (surgery or radiotherapy) is typically made
reactively---after significant growth has already occurred---rather than
proactively based on predicted trajectories. Third, meningioma growth dynamics
are heterogeneous: some tumours remain stable for decades, others exhibit
steady linear growth, and a minority display accelerating or saltatory growth
patterns \cite{hashimoto2012natural,nakamura2003natural}. No established
imaging biomarker reliably predicts which pattern a given tumour will follow.

Quantitative growth prediction from longitudinal MRI could transform
meningioma management by enabling personalised surveillance schedules,
early identification of aggressive growth patterns, and evidence-based
timing of intervention. However, building such predictive models faces a
data bottleneck: longitudinal meningioma MRI datasets are rare, typically
small (tens of patients), and collected retrospectively across heterogeneous
scanners and protocols.

% ---------------------------------------------------------------------------
% 1.2 Technical Challenge
% ---------------------------------------------------------------------------
\subsection{Technical Challenge}
\label{sec:intro:technical}

The central technical challenge is to build a growth prediction model from a
longitudinal cohort of only 33~patients with 2--6~MRI studies each (100~total
observations), while leveraging cross-sectional data from 1,000+ subjects to
learn rich tumour representations. This requires solving three interrelated
problems:

\begin{enumerate}[leftmargin=*]
  \item \textbf{Domain adaptation.} Existing brain MRI foundation models are
    pretrained on glioma data, which differs fundamentally from meningioma in
    tumour morphology, location, and tissue composition. Direct application
    of a glioma-trained encoder to meningioma data produces suboptimal feature
    representations (see \cref{sec:domain_gap}).

  \item \textbf{Structured representation learning.} Growth prediction
    requires features that are not only discriminative but also
    \emph{semantically structured}: tumour volume, spatial location, and
    morphological shape should be disentangled into separate latent
    subspaces amenable to interpretable temporal modelling.

  \item \textbf{Small-sample temporal modelling.} With 33~patients and sparse,
    irregularly-sampled time series, conventional deep learning approaches
    (recurrent networks, Neural ODEs) are catastrophically overparameterised.
    The modelling framework must provide calibrated uncertainty and degrade
    gracefully with limited data.
\end{enumerate}

% ---------------------------------------------------------------------------
% 1.3 Proposed Approach
% ---------------------------------------------------------------------------
\subsection{Proposed Approach}
\label{sec:intro:approach}

We address these challenges through a four-phase pipeline that progressively
transforms raw multi-modal 3D MRI into patient-specific growth predictions
with calibrated uncertainty:

\begin{description}[leftmargin=*,labelwidth=!]
  \item[Phase 1 --- LoRA Adaptation (\cref{sec:lora}).] We adapt
    BrainSegFounder \cite{cox2024brainsegfounder}, a SwinUNETR
    \cite{hatamizadeh2022swinunetr} foundation model pretrained via
    self-supervised learning on 41,400+ brain MRI subjects, to the meningioma
    domain using Low-Rank Adaptation (LoRA) \cite{hu2022lora}. LoRA injects
    trainable low-rank matrices into the encoder's Stages~3--4, adapting
    high-level features while preserving low-level anatomical representations.
    Segmentation on BraTS-MEN~2024 \cite{labella2024brats} serves as the
    proxy task, with optional auxiliary semantic heads providing multi-task
    learning signal.

  \item[Phase 2 --- Supervised Disentangled Projection (\cref{sec:sdp}).] A
    lightweight projection network maps frozen encoder features
    $\mathbf{h} \in \R^{768}$ to a structured latent space
    $\mathbf{z} \in \R^{128}$ partitioned into semantically supervised
    subspaces for volume, location, shape, and a residual component.
    Training combines semantic regression, VICReg-style
    \cite{bardes2022vicreg} variance--covariance regularisation, and distance
    correlation \cite{szekely2007measuring} penalties via a four-phase
    curriculum.

  \item[Phase 3 --- Cohort Encoding (\cref{sec:encoding}).] All longitudinal
    MRI studies from the Andalusian cohort are encoded through the frozen
    pipeline, producing per-patient latent trajectories. ComBat harmonisation
    \cite{johnson2007adjusting,fortin2018harmonization} is assessed and
    applied if significant scanner-induced distributional shift is detected.

  \item[Phase 4 --- Growth Prediction (\cref{sec:growth}).] A three-model
    Gaussian process hierarchy of increasing complexity---Linear Mixed-Effects
    (LME), Hierarchical GP (H-GP), and Partition-Aware Multi-Output GP
    (PA-MOGP)---is fitted on the latent trajectories and evaluated via
    Leave-One-Patient-Out Cross-Validation (LOPO-CV) across 33~folds.
    Volume predictions are decoded through the frozen SDP semantic head and
    compared against observed tumour volumes.
\end{description}

% ---------------------------------------------------------------------------
% 1.4 Contributions
% ---------------------------------------------------------------------------
\subsection{Contributions}
\label{sec:intro:contributions}

This thesis makes the following contributions:

\begin{enumerate}[leftmargin=*]
  \item \textbf{Domain adaptation analysis.} We provide a quantitative
    characterisation of the glioma--meningioma domain gap in the
    BrainSegFounder feature space using Maximum Mean Discrepancy, Centred
    Kernel Alignment, and semantic linear probes (\cref{sec:domain_gap}).

  \item \textbf{LoRA adaptation with semantic supervision.} We demonstrate
    LoRA adaptation of a 3D medical imaging foundation model for meningioma
    segmentation, with ablations across rank, adapter type (LoRA vs.\ DoRA),
    auxiliary semantic heads, and single- vs.\ dual-domain training
    (\cref{sec:lora}).

  \item \textbf{Supervised Disentangled Projection.} We design and
    theoretically justify a non-generative approach to structured latent space
    construction, motivated by the impossibility of unsupervised
    disentanglement \cite{locatello2019challenging} and the failure of
    VAE-based alternatives in our experimental setting (\cref{sec:sdp}).

  \item \textbf{GP hierarchy for latent growth prediction.} We propose a
    principled replacement for Neural ODE-based temporal modelling, with
    closed-form inference, calibrated uncertainty, and a parameter count
    appropriate for the data regime (6--95 parameters vs.\ 3,100+ for the
    Neural ODE) (\cref{sec:growth}).

  \item \textbf{Empirical methodology revision.} We document how empirical
    findings---including shape disentanglement failure ($R^2 \leq 0.11$),
    temporal stationarity of meningioma location, and marginal LoRA benefit
    ($\Delta R^2 = +0.016$)---motivated principled revisions to the pipeline
    architecture (\cref{sec:evaluation}).
\end{enumerate}

% ---------------------------------------------------------------------------
% 1.5 Document Structure
% ---------------------------------------------------------------------------
\subsection{Document Structure}

The remainder of this report is organised as follows.
\Cref{sec:background} reviews the theoretical foundations underlying each
pipeline component. \Cref{sec:data} describes the datasets, preprocessing,
and data splits. \Cref{sec:domain_gap} presents the domain gap analysis
motivating LoRA adaptation. \Cref{sec:lora,sec:sdp,sec:encoding,sec:growth}
detail Phases~1--4 respectively. \Cref{sec:evaluation} consolidates quality
targets, ablation results, and model comparison. \Cref{sec:conclusion}
discusses limitations and future directions.
